\documentclass[11pt]{article}
\usepackage[T1]{fontenc}
\usepackage[utf8]{inputenc}
\usepackage{lmodern,amsmath,amssymb,amsthm,booktabs,array,longtable}
\usepackage[margin=1in]{geometry}
\usepackage[colorlinks=true,urlcolor=blue,citecolor=blue,linkcolor=blue]{hyperref}
\newtheorem{theorem}{Theorem}
\newtheorem{corollary}[theorem]{Corollary}
\newtheorem{lemma}[theorem]{Lemma}
\title{Bounded-error asymptotics for Bern\'ath's conjecture:\\
OEIS A400645, revision 21, COMMENTS,\\
and its rounded-norm siblings}
\author{Haobo Ma\\\small ChronoAI Pte Ltd, Singapore
\and Wenlin Zhang\\\small National University of Singapore\\
\small ChronoAI Pte Ltd, Singapore}
\date{November 9, 2026}
\begin{document}
\maketitle
\begin{abstract}
For Bern\'ath's consecutive-ratio conjecture in OEIS A400645,
revision 21, COMMENTS, we prove the stronger bounded-error asymptotic
$a_n=c\lambda^n+O(1)$ with $c>0$, where $\lambda$ is the largest real
root of $x^8-x^6-x^4-x^2+1$, namely $1.3122819079\ldots$.
The squared terms satisfy a linear recurrence with a controlled rounding
error. A finite seed bounds its entire normalized error tail and ensures
a positive leading coefficient; a three-coordinate contraction then gives
an exponential asymptotic. The same error analysis proves the ratio
assertions for A063827 and A400644, and extends them to nondecreasing
initial data. We also prove the increasing-sequence conjecture in
A289916 and recover its limiting ratio, identified in A289917.
\end{abstract}

\section{The statement and its algebraic constant}
Bern\'ath's sequence \cite{bernath} is defined, with positive indexing, by
\begin{equation}\label{eq:rec}
 a_1=a_2=a_3=0,\quad a_4=1,\qquad
 a_{n+1}=\operatorname{round}\sqrt{s_n},\qquad
 s_n=a_n^2+a_{n-1}^2+a_{n-2}^2-a_{n-3}^2\quad(n\ge4).
\end{equation}
Here rounding means the nearest integer. An integer radicand cannot have
a half-integer square root, so tie conventions are immaterial.
The COMMENTS field of revision 21 says ``Numerical evidence suggests
that the ratio of consecutive terms'' has limit $\sqrt{\mathrm{A289917}}$
and identifies it as the largest real root of
$x^8-x^6-x^4-x^2+1$. We establish that limit,
including existence and positivity of its leading exponential term.

Set
\begin{equation}\label{eq:constant}
 r=\frac{1+\sqrt{13}+\sqrt{2\sqrt{13}-2}}4,
 \qquad \lambda=\sqrt r,\qquad Q(y)=y^4-y^3-y^2-y+1.
\end{equation}
The radical for $r$ is recorded by V\'aclav Kotesovec (2021), and its
largest-root characterization by Linas Vepstas (2024), in A289917
\cite{constant}, revision 48, FORMULA. Both concern the algebraic
constant, rather than the rounded recurrence \eqref{eq:rec}.
For completeness, put $\beta=(1+\sqrt{13})/2$. Then
\[
 Q(y)=(y^2-\beta y+1)(y^2+(\beta-1)y+1).
\]
The first factor has roots $r$ and $r^{-1}$; the second has two distinct
complex-conjugate roots of modulus one, because $0<\beta-1<2$.
Thus $r$ is the largest real root of $Q$, and $\lambda$ is the largest
real root of $Q(x^2)$. The elementary bounds
\begin{equation}\label{eq:rbound}
 \frac{17}{10}<r<\frac74
\end{equation}
follow from $18/5<\sqrt{13}<29/8$ and
$11/5<\sqrt{2\sqrt{13}-2}<19/8$.

\begin{theorem}\label{thm:main}
The recurrence \eqref{eq:rec} defines a unique nonnegative integer
sequence. It is nondecreasing and positive from index four. There is a
constant $c>0$ such that
\[
 a_n\sim c\lambda^n,\qquad
 \lim_{n\to\infty}\frac{a_{n+1}}{a_n}=\lambda.
\]
In particular, the limit asserted in A400645 is the largest real root of
$x^8-x^6-x^4-x^2+1$.
\end{theorem}

\section{Rounding estimates and a positive leading coefficient}
Assume the terms through $a_n$ are nonnegative and nondecreasing. Since
$a_{n-2}\ge a_{n-3}$, one has $s_n\ge a_n^2$. Hence the next term exists
and is at least the integer $a_n$. Induction proves existence,
monotonicity and positivity; uniqueness follows from the deterministic
recurrence and the absence of ties.
Also $a_{n-2}^2\le a_na_{n-1}$, so
\[
 s_n\le (a_n+a_{n-1})^2,\qquad
 a_{n+1}\le a_n+a_{n-1}\quad(n\ge4).
\]
With $q=5/3$, the latter inequality gives, for $m\ge4$ and $j\ge0$,
\begin{equation}\label{eq:growth}
 a_{m+j}\le 2a_m q^j.
\end{equation}
Indeed, the cases $j=0,1$ follow from monotonicity, and the induction
step uses $q^2>q+1$.

Write $t=\sqrt{s_n}$ and $\delta=a_{n+1}-t$, so $|\delta|<1/2$.
The rounding error in the squares satisfies
\begin{equation}\label{eq:error}
 |a_{n+1}^2-s_n|
 \le t+\tfrac14\le 2a_n\qquad(n\ge4).
\end{equation}
Here $s_n\le3a_n^2$, $\sqrt3<7/4$, and $a_n\ge1$ imply the last
inequality. This estimate controls the rounding error by a quantity
growing more slowly than $r^n$.

Shift indices by setting, for $j\ge0$,
\[
 b_j=a_{j+23}^2,\quad
 e_j=b_{j+4}-b_{j+3}-b_{j+2}-b_{j+1}+b_j,
\]
and define the scalar projection
\begin{equation}\label{eq:projection}
 P_j=b_{j+3}+(r-1)b_{j+2}+(r^2-r-1)b_{j+1}-\frac{b_j}{r}.
\end{equation}
The equation $Q(r)=0$ gives the exact identity
\begin{equation}\label{eq:step}
 P_{j+1}=rP_j+e_j.
\end{equation}
Combining \eqref{eq:growth} and \eqref{eq:error} yields
$|e_j|\le4a_{26}q^j$. Therefore the absolutely convergent series
\begin{equation}\label{eq:U}
 U=P_0+\sum_{j=0}^{\infty}\frac{e_j}{r^{j+1}}
   =\lim_{j\to\infty}\frac{P_j}{r^j}
\end{equation}
has the tail bound
\begin{equation}\label{eq:tail}
 |U-P_0|\le\frac{4a_{26}}{r-q}<120a_{26},
\end{equation}
because $r-q>1/30$.

It remains essential to show that $U$ is positive. All coefficients
$r-1$ and $r^2-r-1$ in \eqref{eq:projection} are positive. Monotonicity
gives $b_j\le b_{j+3}$, so
\[
 P_j\ge (1-r^{-1})b_{j+3}>\frac7{17}b_{j+3}.
\]
The finite seed is $a_{26}=324$. To make its exact evaluation transparent,
the initial terms, grouped by their indices, are
\[
\begin{array}{c|rrrrrrrrrrrrr}
 n&1&2&3&4&5&6&7&8&9&10&11&12&13\\\hline
 a_n&0&0&0&1&1&1&2&2&3&4&5&7&9\\[2pt]
 n&14&15&16&17&18&19&20&21&22&23&24&25&26\\\hline
 a_n&12&16&21&28&37&48&63&83&109&143&188&247&324
\end{array}
\]
Every step uses only integer arithmetic: if $k=\lfloor\sqrt{s}\rfloor$,
the rounded value is $k+1$ precisely when $s>k^2+k$.
Consequently \eqref{eq:tail} implies
\begin{equation}\label{eq:positive}
 U>\frac7{17}\,324^2-120\cdot324
   =\frac{73872}{17}>0.
\end{equation}
Thus a finite calculation supplies a quantitative lower bound for the
entire infinite error series, not just numerical evidence of growth.

\section{Convergence of the normalized squares}
Set
\[
 A=\frac{r-1}{r},\quad B=\frac{r^2-r-1}{r^2},\quad C=\frac1{r^4}.
\]
The bounds \eqref{eq:rbound} imply
\begin{equation}\label{eq:ABC}
 0<A<\tfrac12,\qquad 0<B<\tfrac15,\qquad 0<C<\tfrac18.
\end{equation}
For the bound on $B$, use that $1-1/r-1/r^2$ is increasing and at $r=7/4$
equals $5/49<1/5$; the other two bounds are immediate.
Define
\[
 x_j=\frac{b_j}{r^j},\qquad
 L=\frac{U}{r^3(1+A+B-C)}>0,\qquad d_j=x_j-L.
\]
The positivity of the denominator follows from $1+A+B-C>7/8$.
Dividing \eqref{eq:projection} by $r^{j+3}$ gives
\[
 d_{j+3}=f_j-Ad_{j+2}-Bd_{j+1}+Cd_j,\qquad
 f_j=\frac{P_j}{r^{j+3}}-\frac U{r^3}\longrightarrow0.
\]
For the weighted absolute sum
\[
 M_j=|d_{j+2}|+\frac25|d_{j+1}|+\frac3{20}|d_j|,
\]
the triangle inequality and \eqref{eq:ABC} give
\begin{equation}\label{eq:contraction}
 M_{j+1}\le\frac9{10}M_j+|f_j|.
\end{equation}
Indeed, the coefficients of $|d_{j+2}|$, $|d_{j+1}|$ and $|d_j|$
are at most $9/10$, $7/20\le(9/10)(2/5)$ and
$1/8\le(9/10)(3/20)$, respectively.
To see that \eqref{eq:contraction} implies $M_j\to0$, fix
$\varepsilon>0$ and choose $J$ with $|f_j|\le\varepsilon/10$ for $j\ge J$.
Iteration gives
$M_{J+k}\le(9/10)^kM_J+\varepsilon$; taking a limit superior and then
letting $\varepsilon\downarrow0$ proves the assertion.
Hence $d_j\to0$, and
\[
 \frac{a_{j+23}^2}{r^j}\longrightarrow L>0.
\]
Taking nonnegative square roots proves
$a_n\sim c\lambda^n$ with $c=\sqrt L\,\lambda^{-23}>0$.
More explicitly, division of successive normalized squares gives
$a_{n+1}^2/a_n^2\to r$; positivity then gives
$a_{n+1}/a_n\to\sqrt r$. This completes the proof of
Theorem~\ref{thm:main}.

\section{Bounded error and positive-norm siblings}
\begin{lemma}\label{lem:forcing}
Let a monic real polynomial $p$ have distinct roots, a root $R>1$,
and all its other roots of modulus at most one. If a real sequence
obeys the recurrence with characteristic polynomial $p$, with additive
error $e_j=O(q^j)$ for some $1<q<R$, then
\[
 b_j=L R^j+O(q^j)
\]
for some real $L$. If $b_j\sim L R^j$ with $L>0$ and the error also
satisfies $e_j=O(b_j^{1/2})$, then
$b_j=L R^j+O(R^{j/2})$.
\end{lemma}
\begin{proof}
The companion matrix is diagonalizable over $\mathbb C$.
Its eigenvectors form a Vandermonde matrix, whose determinant is a
product of nonzero root differences. In these coordinates each
component satisfies $v_{j+1}=\zeta v_j+\kappa e_j$, whence
\[
 v_j=\zeta^jv_0+\kappa\sum_{k<j}\zeta^{j-1-k}e_k.
\]
For $|\zeta|\le1$ the sum is $O(q^j)$. For $\zeta=R$, the series
$\sum_{k\ge0}e_k/R^{k+1}$ converges; replacing the finite sum by that
series leaves a tail of size $O(q^j)$. Transforming back proves the
first assertion. Under the additional hypotheses $e_j=O(R^{j/2})$,
the same argument applies with $q=\sqrt R$. The resulting leading
coefficient is the same $L$, since division by $R^j$ determines it
uniquely.
\end{proof}

\begin{corollary}\label{cor:bounded}
For the sequence in Theorem~\ref{thm:main},
$a_n=c\lambda^n+O(1)$, and
$a_{n+1}/a_n=\lambda+O(\lambda^{-n})$.
\end{corollary}
\begin{proof}
Theorem~\ref{thm:main} and \eqref{eq:error} give
$e_j=O(r^{j/2})$ for its shifted squared sequence.
The factorization of $Q$ and Lemma~\ref{lem:forcing} give
$a_n^2=c^2r^n+O(r^{n/2})$. Rationalizing the square-root difference
gives $|a_n-c\lambda^n|=O(1)$, since its denominator is at least
$c\lambda^n$. Division of consecutive terms gives the ratio estimate.
\end{proof}

\begin{corollary}\label{cor:large}
Replace the initial tuple in \eqref{eq:rec} by any nondecreasing
nonnegative integer tuple with $a_4\ge292$. The resulting signed
recurrence satisfies $a_n=c\lambda^n+O(1)$ for some $c>0$.
\end{corollary}
\begin{proof}
Existence, monotonicity and \eqref{eq:growth}--\eqref{eq:error} hold
unchanged. Set $b_j=a_{j+1}^2$ and $A=a_4$. The projection
\eqref{eq:projection} has $P_0>(7/17)A^2$, while its entire error
tail is less than $120A$. Since
$(7/17)292^2-120\cdot292=1168/17>0$, the normalized projection has
positive limit for every $A\ge292$. The contraction proof and
Corollary~\ref{cor:bounded} now apply with this index shift.
\end{proof}

\begin{corollary}\label{cor:norm}
Let $m=2$ or $3$ and start with nondecreasing nonnegative integer data
$x_0,\ldots,x_{m-1}$. Define
\[
 x_{n+m}=\operatorname{round}\sqrt{\sum_{i=0}^{m-1}x_{n+i}^2}.
\]
For $m=2$ assume $x_1\ge2$; for $m=3$ assume $x_2\ge1$.
If $R_m$ is the positive root above one of
$p_m(y)=y^m-y^{m-1}-\cdots-y-1$, then for some $C>0$,
\[
 x_n=C R_m^{n/2}+O(1),\qquad
 \frac{x_{n+1}}{x_n}=\sqrt{R_m}+O(R_m^{-n/2}).
\]
In particular, the asserted limits in A063827, revision 12, COMMENTS,
and A400644, revision 28, COMMENTS, hold for their specified initial
data: $\sqrt{(1+\sqrt5)/2}$ and the square root of the tribonacci
constant, respectively.
\end{corollary}
\begin{proof}
The sequence exists uniquely and remains nondecreasing, since the
radicand is at least the square of the last term. If it were bounded,
it would eventually be a constant integer $K$. The hypotheses imply
$(\sqrt m-1)K>1/2$, contradicting
$\operatorname{round}(\sqrt m K)=K$. Thus $x_n\to\infty$.

For $m=2$ the roots are $R_2=(1+\sqrt5)/2>\sqrt2$ and $-R_2^{-1}$.
For $m=3$, the cubic has just one real root $R_3>\sqrt3$:
its extrema at $-1/3$ and $1$ both have negative values, and
$p_3(\sqrt3)=2\sqrt3-4<0$. The other two roots are conjugate and
have modulus $R_3^{-1/2}<1$, by their product. All these roots are
distinct. Choose $\sqrt m<q<R_m$ and a late index $N$ such that
$\sqrt m+1/(2x_N)\le q$ and
$A=x_{N+m-1}>2/(R_m-q)$. Then
$x_{N+m-1+j}\le A q^j$ by induction.

Set $b_j=x_{N+j}^2$ and
$e_j=b_{j+m}-\sum_{i<m}b_{j+i}$.
The rounding bound gives
$|e_j|\le\sqrt m\,x_{N+j+m-1}+1/4\le2Aq^j$.
Writing $R=R_m$, use the positive-coefficient projection
\[
 P_j=\begin{cases}
 b_{j+1}+(R-1)b_j,&m=2,\\
 b_{j+2}+(R-1)b_{j+1}+R^{-1}b_j,&m=3.
 \end{cases}
\]
Direct substitution of $p_m(R)=0$ gives $P_{j+1}=RP_j+e_j$.
Therefore
\[
 U=\lim_j P_j/R^j
 \ge A^2-\frac{2A}{R-q}>0.
\]
Lemma~\ref{lem:forcing} gives $b_j=L R^j+O(q^j)$, so
$U=L p_m'(R)$. The derivative is positive at the largest real root,
hence $L>0$. Now $e_j=O(b_{j+m-1}^{1/2})=O(R^{j/2})$.
The second part of the lemma and rationalization give the asserted
bounded error and ratio estimate. A063827 starts at $(1,2)$;
A400644 starts at $(0,0,1)$ with positive indexing, so both satisfy
the hypotheses.
\end{proof}

\section{Algebraic companions and source coverage}
\begin{corollary}
For every fixed integer $h\ge0$, $a_{n+h}/a_n\to\lambda^h$, and
$\log(a_n)/n\to\log\lambda$.
\end{corollary}
\begin{proof}
Both statements follow directly from $a_n\sim c\lambda^n$ with $c>0$;
the logarithm is defined for $n\ge4$.
\end{proof}

The unrounded companion in A116732 \cite{linear} has
$u_0=u_1=u_2=0$, $u_3=1$, and
$u_n=u_{n-1}+u_{n-2}+u_{n-3}-u_{n-4}$ for $n\ge4$.
Its generating function is $z^3/Q(z)$, as recorded in its FORMULA field.
The four simple roots of $Q$ described above give, by partial fractions,
\[
 u_n=\frac{r^n}{Q'(r)}+O(1),\qquad
 \frac{u_{n+1}}{u_n}\longrightarrow r=\lambda^2.
\]
For the coefficient, at $z=1/r$ the derivative of $Q(z)=z^4Q(1/z)$
is $-Q'(r)/r^2$, so the coefficient of $(1-rz)^{-1}$ in $z^3/Q(z)$
is $1/Q'(r)$. The other poles have moduli $1,1,r$, giving bounded
coefficients. Moreover
$Q'(r)=(r-r^{-1})(r^2+(\beta-1)r+1)>0$.
This explains the square substitution in A400645 without identifying
its rounded squares with the distinct sequence $u_n$.

\begin{corollary}\label{cor:reverse}
The unique extension of $u_n$ to all integer indices satisfies
$u_n=-u_{2-n}$, the identity recorded by Michael Somos (July 25, 2025)
in A116732, revision 33, FORMULA.
\end{corollary}
\begin{proof}
The recurrence has invertible last coefficient, so determines the
extension uniquely, with $u_{-1}=-1$. Its palindromic characteristic
polynomial makes $-u_{2-n}$ another solution; its values at $n=0,1,2,3$
are $0,0,0,1$. Uniqueness proves the identity.
\end{proof}

\begin{corollary}\label{cor:kimberling}
The increasing-sequence conjecture in Kimberling's A289916, revision 11,
COMMENTS, holds: its coefficients $v_n$ are strictly increasing for
$n\ge0$. Moreover $v_n=D r^n+O(1)$ for some $D>0$, so their
consecutive-ratio limit is the number defined in A289917 \cite{kimberling}.
\end{corollary}
\begin{proof}
Colin Barker's July 19, 2017 formula in A289916 gives
\[
 F(z)=\sum_{n\ge0}v_nz^n
 =\frac{(1+z)^2(1-z+z^2-z^3+z^4-z^5+z^6)}
 {(1-z+z^2)Q(z)}
 =\frac{(1+z)^2(1+z^7)}{(1+z^3)Q(z)}.
\]
For completeness, this formula follows directly from its NAME:
if $w_k=\operatorname{round}(9(k+1)/7)$, then $w_{k+7}=w_k+9$ and
$(w_0,\ldots,w_6)=(1,3,4,5,6,8,9)$.
Consequently
\[
 \sum_{k\ge0}w_kt^k
 =\frac{1+2t+t^2+t^3+t^4+2t^5+t^6}{(1-t)(1-t^7)}.
\]
Substitute $t=-z$ and invert; cancellation gives $F$ above.
Now let $H(z)=1/((1+z^3)Q(z))=\sum h_nz^n$, with $h_n=0$ for $n<0$.
Since
\[
 (1+z^3)Q(z)=1-z-z^2-z^5-z^6+z^7,
\]
one has $h_0=1$ and, for $n\ge1$,
\[
 h_n-h_{n-1}=h_{n-2}+h_{n-5}+h_{n-6}-h_{n-7}.
\]
Induction shows that $h_n$ is nondecreasing, and strictly increasing
for $n\ge2$: the right side is nonnegative and contains
$h_{n-2}\ge1$ then. Multiplication by $(1+z)^2(1+z^7)$ preserves
nonnegative differences; for $n=1$, $v_1-v_0=2$, and for $n\ge2$
the $h_n$ summand gives strictness. This proves the conjecture,
whose exact wording is ``Conjecture: the sequence is strictly increasing.''

The first expression for $F$ has simple denominator roots: those of
$Q$ and the two roots of $1-z+z^2$; these sets are disjoint since
$Q(\zeta)=3(1-\zeta)\ne0$ for $\zeta^2-\zeta+1=0$.
At $z=1/r$ its numerator and the quadratic factor are positive and
$Q'(1/r)=-Q'(r)/r^2<0$. Thus the coefficient of $(1-rz)^{-1}$
is positive. The other poles have modulus at least one, so partial
fractions give $v_n=D r^n+O(1)$.
\end{proof}

\begin{corollary}\label{cor:descartes}
For real $y\ne0$,
\[
 \left(\sum_{i=0}^3y^{-i}\right)^2-3\sum_{i=0}^3y^{-2i}
 =-\frac{2(y^2+1)Q(y)}{y^6}.
\]
Consequently the largest real solution of the left side equal to zero
is $r$. This proves the algebraic identity recorded by Charles L. Hohn
(November 13, 2025), A289917, revision 48, FORMULA.
\end{corollary}
\begin{proof}
Multiplication by $y^6$ and expansion give the displayed identity;
$y^2+1>0$ and the factorization of $Q$ give the root assertion.
\end{proof}

\subsection*{Coverage of the posing source and its companions}
All ``Proved'' labels below refer to statements actually asserted or
conjectured in the indicated source. Further statements are identified
as extensions or as not asserted.
{\small
\begin{longtable}{@{}p{.29\linewidth}p{.65\linewidth}@{}}
\toprule
Source item & Mathematical coverage\\\midrule
A400645, rev. 21, NAME / Mathematica / COMMENTS &
Proved: suggested ratio limit for the specified initial data,
Theorem~\ref{thm:main}. Not posed; proved: bounded absolute error and
exponential ratio convergence, Corollary~\ref{cor:bounded}.\\
A063827, rev. 12, COMMENTS &
Proved: the asserted square-root golden-ratio limit,
Corollary~\ref{cor:norm}. Its initial data are $(1,2)$ \cite{two}.\\
A400644, rev. 28, COMMENTS &
Proved: the suggested square-root tribonacci limit,
Corollary~\ref{cor:norm}. Its positive-index initial data are $(0,0,1)$
\cite{three}.\\
Nondecreasing initial-data extension for $m=2,3$ &
Not posed; proved under the explicit size hypotheses of
Corollary~\ref{cor:norm}. The all-one two-term sequence is stationary.\\
Large nondecreasing data for the signed rule &
Not posed; proved when the last of the four initial terms is at least
$292$, Corollary~\ref{cor:large}.\\
A289917, rev. 48, FORMULA &
Prior identities cited and proved: Kotesovec's radical (2021),
Vepstas's quartic root (2024), and Hohn's sum identity (2025);
\eqref{eq:constant} and Corollary~\ref{cor:descartes}.\\
A289916, rev. 11, COMMENTS / FORMULA &
Proved: Kimberling's strictly increasing conjecture.
Barker's rational generating function (2017) is cited and derived;
Corollary~\ref{cor:kimberling} also proves the ratio used in
A289917's NAME.\\
A116732, rev. 33, NAME / FORMULA &
Prior recurrence and generating function cited; positive asymptotic
proved above. Somos's reversal identity (2025) is cited and reproved
in Corollary~\ref{cor:reverse}.\\
A116732, rev. 33, COMMENTS / EXAMPLE &
Not covered: unspecified ``expected divisibility properties'' and
the partially ordered partition interpretation.\\
A289917, rev. 48, COMMENTS; other geometric cross-references &
Not covered: kissing-sphere geometry, radius interpretation and
the related geometric configurations.\\
Arbitrary positive data for the signed four-term rule &
Not asserted in the source; fails for the all-one sequence, whose
ratio is one. This is not a refutation of A400645.\\
\bottomrule
\end{longtable}}
The initial conditions matter: the constant sequence $a_n=1$ satisfies
the same four-term rule, since $\operatorname{round}\sqrt2=1$, and has
ratio one. Thus Theorem~\ref{thm:main} is not a statement about every
positive initial tuple.

\paragraph{Verification coverage.}
The pinned Lean file
\href{https://github.com/the-omega-institute/trureturing/blob/6974dba01daa1a25cc74b87a86c191e95fe2c681/D5/S1/Recurrence/OEISA400645.lean}{\texttt{D5/S1/Recurrence/OEISA400645.lean}}
contains the declaration
\texttt{D5.S1.Recurrence.OEISA400645.result} with statement
\texttt{claim}: existence for \texttt{IsSourceSequence}, the
consecutive-ratio limit for every sequence satisfying it, and the largest
real-root characterization. The positive asymptotic appears in its
private lemma \texttt{square\_asymptotic}. The arguments above supply
written proofs of the bounded-error refinement, fixed-lag limits,
general initial-data extension, and algebraic companion corollaries.
The second pinned Lean file,
\href{https://github.com/the-omega-institute/trureturing/blob/6974dba01daa1a25cc74b87a86c191e95fe2c681/D5/S1/Recurrence/OEISA400644.lean}{\texttt{D5/S1/Recurrence/OEISA400644.lean}},
has declaration \texttt{D5.S1.Recurrence.OEISA400644.result},
with statement \texttt{targetClaim}:
existence and uniqueness for its positive-index initial data, positivity,
the infinite consecutive-ratio limit, and its largest-positive-sextic-root
characterization. It does not state the general initial-data or
bounded-error extensions. The two-term and A289916 statements have written
proofs here. The accompanying integer-arithmetic checker tests 512 terms
of each rounded sequence, the displayed source terms, the main seed and
rational constants, the companion identities and monotonicity; those
finite tests are distinct from the infinite proofs.

\begin{thebibliography}{9}
\bibitem{bernath} B. Bern\'ath.
OEIS A400645, revision 21, NAME, COMMENTS, Mathematica and CROSSREFS
fields; authored October 2, 2026.
\url{https://oeis.org/A400645}.
\bibitem{constant} C. Kimberling.
OEIS A289917, revision 48, FORMULA field, with contributions by
V. Kotesovec (August 27, 2021) and L. Vepstas (February 6, 2024).
\url{https://oeis.org/A289917}.
\bibitem{linear} R. K. Guy.
OEIS A116732, revision 33, NAME and FORMULA fields.
\url{https://oeis.org/A116732}.
\bibitem{two} H. Bottomley.
OEIS A063827, revision 12, NAME and COMMENTS fields.
\url{https://oeis.org/A063827}.
\bibitem{three} B. Bern\'ath.
OEIS A400644, revision 28, NAME and COMMENTS fields.
\url{https://oeis.org/A400644}.
\bibitem{kimberling} C. Kimberling.
OEIS A289916, revision 11, NAME, COMMENTS and FORMULA fields;
authored July 18, 2017, with formulas by C. Barker (July 19, 2017).
\url{https://oeis.org/A289916}.
\end{thebibliography}
\medskip
\noindent\small Licensed under Creative Commons Attribution 4.0
International (CC BY 4.0).
\end{document}
